\section{Ontological Limits}

\subsection{Morphogenesis as the Stable Attractor}

Reviewing the progression from Phases 1 through the speculative extensions, we arrive at what we consider the paper's central theoretical claim. If the pattern of bottleneck migration is genuine (an empirical question), and if it persists at higher levels of architectural complexity (a conjecture), then the system's long-term trajectory is not toward any particular fixed architecture but toward a sustained capacity for architectural change. The cycle of detecting a bottleneck, fissioning a role, and observing the emergence of a new bottleneck would repeat at every scale: within individual agent roles, across pipelines, across populations, and at the meta-architectural level.

This characterization aligns with Prigogine and Nicolis' \cite{prigogine1977} theory of dissipative structures: systems maintained far from thermodynamic equilibrium through continuous flows of energy or information. The agentic ecosystem, on this view, persists not despite its continuous reorganization but because of it. If bottlenecks cease to migrate, the system has reached equilibrium, which in this context would indicate stagnation. The connection to Friston's \cite{friston2010} Free Energy Principle is also suggestive: the Architect's self-model functions as a generative model of the system, bottleneck detection corresponds to surprise minimization, and role fission corresponds to active inference. We stress, however, that these are analogies, not demonstrations; whether the mathematical structure of dissipative systems or free energy minimization actually applies to agentic architectures is a question for future formal analysis.

\subsection{Self-Reference and Incompleteness}

If a system constructs an explicit model of its own capabilities and limitations (as we hypothesize for Phase 8 and beyond), it enters a self-referential regime that can be analyzed using two classical results. Lawvere's fixed-point theorem \cite{lawvere1969} guarantees that in any cartesian closed category with a point-surjective morphism $\varphi: A \to B^A$, every endomorphism $f: B \to B$ has a fixed point. A verifier that verifies itself is precisely such a fixed point. Lawvere's theorem establishes that such fixed points exist in any sufficiently rich formal setting; the question is not whether self-referential equilibria are possible but what properties they have. G\"odel's first incompleteness theorem \cite{godel1931} constrains the answer: any self-model powerful enough to be useful is necessarily incomplete. The system can know that it has blind spots but cannot enumerate all of them from within.

Taken together, these results suggest that self-reconfiguring agentic systems can achieve stable but irreducibly partial self-understanding. This connects to Hofstadter's \cite{hofstadter1979} theory of strange loops, Yanofsky's \cite{yanofsky2003} unified treatment of self-referential paradoxes, and Barwise and Moss's \cite{barwise1996} non-well-founded set theory. Von Foerster's \cite{vonfoerster2003} second-order cybernetics arrives at the same structure from the operational side: when an observer recursively observes its own observations, the process converges to stable \emph{eigenbehaviors}---fixed points that the system can reliably maintain, not because they have been externally certified but because they are self-consistent solutions of the recursive process. Lawvere's theorem guarantees the existence of such fixed points in any sufficiently rich formal setting; Von Foerster's framework explains how they are realized in practice through recursive interaction rather than top-down computation. We note that applying these mathematical results to engineered software systems involves a degree of analogical reasoning that may or may not survive formal scrutiny; we present them as a conceptual framework, not as a proof.

\subsection{Ontological Pluralism and the Observer}

Finally, we note three meta-level observations that bear on the framework's interpretation. First, a single system event (e.g., a role fission) can be described as ``role differentiation'' in organizational vocabulary, ``symmetry breaking'' in thermodynamic vocabulary \cite{anderson1972}, ``strategy differentiation'' in game-theoretic vocabulary \cite{maynard-smith1982}, or ``niche divergence'' in ecological vocabulary \cite{hutchinson1957}. Each description language reveals structures that the others occlude; a system described in only one vocabulary may exhibit systematic blind spots undetectable from within that vocabulary.

Second, if all components of a system are replaceable (today's LLM agent could be tomorrow's human engineer), the system's identity resides not in its substrate but in its relational structure: its information flow topology, preference constraints, and coordination protocols. This is the category-theoretic perspective on identity \cite{mac_lane1971}, and it suggests that agentic system design should prioritize protocol-level interoperability over implementation-level homogeneity.

Third, the act of describing an agentic system is itself a component of the system it describes: a researcher analyzing the Wallfacer experiment functions as a Strategist (generating hypotheses), the analysis serves as an Executor (producing formalized descriptions), and the resulting paper acts as a Cartographer (creating artifacts that inform future experimental design). This self-referential observation, anticipated by Maturana and Varela's \cite{maturana1980} concept of autopoiesis, does not invalidate the framework but does caution against treating any description of an evolving system as final.
